\label{ch:summary}

This chapter concludes the thesis by summarizing the contributions of the proposed thermal infrared fall detection system (Section~\ref{sec:summary}) and outlining directions for future research (Section~\ref{sec:outlook}).

%% ====================================================================
\section{Summary}
\label{sec:summary}
%% ====================================================================

Falls represent one of the most significant health hazards facing the elderly population worldwide. According to the World Health Organization, approximately 684{,}000 fatal falls occur globally each year, and adults aged 60 and older account for the largest share of these incidents. In residential care facilities, timely detection of falls can dramatically reduce the severity of outcomes---studies have shown that lying on the floor undetected for more than one hour is independently associated with serious injury and increased mortality, even when the initial fall itself is not severe. Automated fall detection systems based on computer vision offer a promising approach to continuous, non-intrusive monitoring. However, conventional RGB camera-based methods raise fundamental privacy concerns: residents, families, and regulatory bodies are often reluctant to accept continuous video surveillance that captures identifiable imagery of individuals in their private living spaces. Moreover, RGB systems are inherently unreliable under the low-light conditions common in residential rooms during nighttime---precisely the period when fall risk is elevated due to disorientation, medication effects, and reduced visibility. Thermal infrared (IR) imaging provides a compelling alternative modality. By capturing emitted body heat rather than reflected visible light, thermal cameras preserve occupant privacy (no identifiable facial features or clothing details are recorded) and operate robustly regardless of ambient illumination. The central technical challenge, however, is that thermal IR images---particularly from low-cost thermopile sensor arrays such as the Heimann HTPA series used in the SensIR project---exhibit substantially lower spatial resolution and thermal contrast compared to RGB, rendering fine-grained human pose estimation and temporally nuanced action recognition considerably more difficult.

To address these challenges, this thesis presented an end-to-end multi-task framework for real-time fall detection from thermal infrared imagery. The system is built upon the RTMDet~\cite{lyu2022rtmdet} anchor-free detection architecture and integrates three tightly coupled tasks within a unified pipeline. In the first stage, a person detection head localizes human subjects with bounding box regression using Quality Focal Loss~\cite{li2020generalized} and GIoU loss~\cite{rezatofighi2019generalized}, while a pose estimation head simultaneously predicts seven anatomical keypoints (head, shoulder center, both hands, hip center, both feet) via SimDR~\cite{li20212d}---a disentangled 1D coordinate classification method that decomposes the 2D localization problem into two independent 1D classification tasks along the $x$- and $y$-axes. In the second stage, a temporal action classification head built upon a Gated Recurrent Unit (GRU)~\cite{cho2014learning} processes sequences of normalized keypoint features extracted from tracked individuals. The proposed dual-task design of this action head simultaneously performs four-class posture classification (standing, walking, sitting, lying) and binary fall detection, sharing a common GRU temporal encoder so that posture understanding directly informs fall discrimination. The backbone employs MobileViT-S~\cite{mehta2022mobilevit}, a hybrid CNN-Transformer architecture with only 4.9M parameters that combines the local feature extraction strength of convolutional layers with the global contextual modeling capability of self-attention, paired with a CSPNeXtPAFPN~\cite{lyu2022rtmdet} neck for multi-scale feature aggregation across strides $\{8, 16, 32\}$. RoI features for each detected person are extracted via RoIAlign~\cite{he2017mask} at a fixed output size of $48 \times 48$ pixels. The two-stage training strategy---where Stage~1 jointly optimizes detection and pose estimation, and Stage~2 freezes all Stage~1 parameters before training only the randomly initialized temporal action head---enables independent development, evaluation, and modular replacement of each component without retraining the entire pipeline. Multi-object tracking via ByteTrack~\cite{zhang2022bytetrack} with Kalman filter motion prediction maintains consistent person identities across frames, enabling the GRU to process coherent per-person keypoint sequences of up to 30 frames.

Comprehensive experimental evaluation on the SensIR thermal IR dataset yielded four principal findings. First, the choice of pose estimation representation proved to be the single most impactful design decision in the entire pipeline. As reported in Table~\ref{tab:pose_comparison}, SimDR achieved a COCO Keypoint AP of 0.802, outperforming the 2D Heatmap baseline (0.573) and the Direct Coordinate Regression approach (0.570) by a relative margin of approximately 40\%. At the stricter AP@0.75 threshold, SimDR's advantage was even more pronounced (0.874 vs.\ 0.574 and 0.595). The disentangled 1D classification formulation with KL divergence loss against soft Gaussian targets was particularly effective for the small RoI sizes ($48 \times 48$) characteristic of thermal IR images, where the quantization error inherent in 2D heatmap argmax decoding is proportionally large relative to the spatial resolution. This finding has broader implications: SimDR's advantages may extend to other low-resolution pose estimation tasks beyond thermal imaging.

Second, the ablation study on input feature representations (Table~\ref{tab:feature_comparison}) demonstrated that raw keypoint features---encoding absolute position ($x$, $y$), visibility, and inter-frame displacement ($\Delta x$, $\Delta y$) in 35 dimensions---consistently outperformed more complex alternatives. Bone skeleton features (36 dimensions), which encode translation-invariant relative joint displacements, achieved a Fall F1 of only 0.680, while adding 96-dimensional RoI appearance features further degraded performance to 0.641 due to the low discriminative value of thermal IR appearance and the increased overfitting risk from the $3.7\times$ dimensionality expansion. The decisive advantage of raw keypoint features confirms that absolute spatial information---particularly the vertical position of the head relative to the bounding box---is the most discriminative cue for fall detection, a finding that is intuitive given that falls fundamentally involve a transition from upright to horizontal body configuration.

Third, the fall classification strategy comparison (Table~\ref{tab:fall_comparison}) revealed a nuanced trade-off between architectural completeness and data efficiency. The proposed dual-task method achieved a Fall F1 of 0.669 with \emph{perfect precision} (1.000)---producing zero false positives on the validation set---while simultaneously providing clinically valuable four-class posture classification output. The simpler Binary-BCE baseline achieved a higher overall Fall F1 of 0.748 (Table~\ref{tab:loss_comparison}), driven by superior recall (0.629 vs.\ 0.502). This performance gap was attributed to the dual-task model's higher representational capacity leading to underfitting on the limited training set of 4{,}838 instances: the model must simultaneously learn to discriminate four posture classes and detect falls, distributing its learning capacity across two objectives. The loss function ablation further confirmed that on small datasets, simpler optimization landscapes (plain BCE with class weighting) outperform more sophisticated alternatives (focal loss, hard negative mining), which represent architectural investments that pay off only with sufficient training data. With synthetic data augmentation, the proposed dual-task method improves to a Fall F1 of 0.776 (recall: 0.655, precision: 0.951), surpassing the Binary-BCE baseline and confirming that the dual-task architecture scales favorably with additional data.

Fourth, computational profiling (Table~\ref{tab:computational_cost}) confirmed that the complete system is highly efficient: 9.3M total parameters with the full inference pipeline achieving approximately 80~FPS on an NVIDIA TITAN RTX GPU. The temporal action head contributes negligible overhead---only 31{,}617 parameters (0.3\% of the total model) and 0.8~ms per frame---demonstrating that the GRU-based skeleton classifier adds rich temporal reasoning at virtually no computational cost. ByteTrack multi-object tracking adds less than 0.5~ms per frame. These characteristics make the system well-suited for deployment on edge devices equipped with FPGA or NPU acceleration, aligning with the SensIR project's target of $>$20~FPS on embedded hardware.

Despite the Binary-BCE baseline achieving higher Fall F1 on the current dataset, the dual-task architecture was deliberately selected as the final system because it provides the most architecturally complete pipeline for real-world clinical deployment. This decision was motivated by four considerations. First, the four-class posture classification output is directly useful for clinical monitoring applications that extend beyond acute fall detection---for instance, identifying prolonged lying episodes that may indicate a post-fall inability to self-recover, tracking declining daily mobility as an early warning of deteriorating health, or monitoring sleep patterns in residential care. Second, the perfect fall precision (1.000) guarantees zero false alarms, a property that is critical for clinical trust and adoption. Healthcare staff subject to frequent false alerts develop alarm fatigue, eventually disregarding or delaying response to notifications~\cite{sendelbach2013alarm}; a system with zero false positives avoids this entirely. Third, the shared GRU encoder that jointly learns posture representations and fall discrimination forms a closed-loop design: understanding that a person is in a lying posture directly constrains the fall detection decision, and distinguishing between controlled lying-down transitions and uncontrolled falls requires precisely the posture-level understanding that the auxiliary task provides. Fourth, the architecture's higher representational capacity scales favorably with additional training data, as confirmed by the synthetic data augmentation results: the proposed method achieves Fall F1 of 0.776 with synthetic augmentation, surpassing both its baseline (0.669) and the Binary-BCE alternative (0.748).


In conclusion, this thesis demonstrates that thermal infrared fall detection can achieve real-time, privacy-preserving performance through a lightweight multi-task architecture that jointly addresses person detection, pose estimation, and temporal action classification within a single framework. The combination of SimDR-based pose estimation---which dramatically outperforms conventional heatmap and regression alternatives on low-resolution thermal features---with a dual-task GRU temporal classifier provides both quantitative accuracy and clinical interpretability. The system's compact model size of 9.3M parameters, modular two-stage training pipeline, and approximately 80~FPS throughput make it a practical foundation for autonomous fall detection in residential and clinical care facilities, where privacy preservation, real-time responsiveness, and deployment on resource-constrained edge hardware are essential requirements.

%% ====================================================================
\section{Outlook}
\label{sec:outlook}
%% ====================================================================

While the proposed system establishes a functional end-to-end pipeline for thermal IR fall detection and demonstrates the viability of the multi-task approach, several directions for future research can address its current limitations and substantially extend its capabilities. This section discusses five such directions, ordered roughly by their expected impact and feasibility.

\subsection{Synthetic Data Augmentation and Dataset Expansion}
\label{sec:outlook_synth}

The most pressing limitation of the current system is the modest size of the training dataset. The SensIR dataset contains only 4{,}838 training instances with 1{,}482 falling examples---a volume that is insufficient to fully exploit the dual-task model's representational capacity and that contributes directly to the performance gap relative to simpler single-task baselines. As demonstrated by the loss function ablation (Table~\ref{tab:loss_comparison}), increasing model complexity or adding sophisticated loss components (focal loss, hard negative mining) yielded diminishing or even negative returns on this data volume, indicating that the primary bottleneck is data quantity rather than architectural expressiveness.

Two complementary strategies for synthetic data generation will be investigated. The first involves physics-based simulation of thermal IR sequences using parametric human body models (e.g., SMPL~\cite{loper2015smpl}) rendered under simulated thermal emission profiles. By varying body proportions, clothing insulation properties, ambient temperature, and camera viewpoint, a diverse corpus of synthetic falling and non-falling sequences can be generated with automatic ground-truth annotations for bounding boxes, keypoints, and action labels. The second strategy leverages domain adaptation from existing large-scale RGB fall detection datasets---such as the UR Fall Detection dataset~\cite{kwolek2014human}, the Le2i Fall Detection dataset~\cite{charfi2013optimized}, or the NTU RGB+D dataset~\cite{shahroudy2016ntu}---by learning an image-to-image translation mapping from RGB to thermal IR appearance using generative adversarial networks (GANs) or diffusion models, while preserving the original pose and action annotations. This approach benefits from the substantially larger volume and diversity of existing RGB fall datasets without requiring additional physical data collection.

The proposed dual-task architecture is expected to benefit disproportionately from additional training data compared to simpler baselines. The Binary-BCE model, with its single optimization objective and minimal architectural complexity, likely approaches its performance ceiling on the current dataset; additional data would provide diminishing returns. In contrast, the dual-task model's shared GRU encoder, which must learn both posture discrimination and fall detection simultaneously, requires a richer and more diverse training signal to fully exploit the representational synergies between these two tasks. Preliminary experiments with synthetic data augmentation confirm this hypothesis: the proposed method achieves a Fall F1 of 0.776 with synthetic augmentation, surpassing both the base proposed result (0.669) and the Binary-BCE baseline (0.748).

\subsection{Advanced Temporal Modeling}
\label{sec:outlook_temporal}

The current temporal action head employs a single-layer GRU processing fixed-length keypoint sequences of $T = 8$ frames with a per-frame MLP encoder that independently projects each frame's 35-dimensional keypoint vector into a 64-dimensional latent space before temporal aggregation. While this design is computationally efficient and sufficient for the current dataset, it has several inherent limitations that more expressive temporal architectures could address.

First, the fixed sequence length of $T = 8$ frames constrains the temporal receptive field to approximately 0.8 seconds at 10~FPS. While this window captures the acute phase of most falls, it is insufficient for distinguishing slow, controlled transitions (e.g., gradually sitting down over 3--5 seconds) from rapid falls that unfold over a similar duration but with different acceleration profiles. A Transformer-based temporal encoder~\cite{vaswani2017attention} with multi-head self-attention and sinusoidal positional encoding would naturally accommodate variable-length input sequences and enable the model to attend selectively to the most discriminative frames within arbitrarily long temporal windows. The attention weight distribution would also provide interpretability---revealing which temporal moments the model considers most diagnostic of a fall.

Second, the per-frame MLP encoder treats the seven keypoints as a flat feature vector, discarding the spatial topology of the human body. Spatial-Temporal Graph Convolutional Networks (ST-GCN)~\cite{yan2018spatial} and their successors~\cite{shi2019two,chen2021channel} operate directly on the skeleton graph structure, applying graph convolutions that aggregate information along anatomical connections (e.g., head$\to$shoulder$\to$hip$\to$foot). For our seven-keypoint skeleton, the graph structure is admittedly simple, but explicit graph reasoning could still improve the model's ability to detect coordinated multi-joint patterns characteristic of falls---such as simultaneous head drop and leg buckling---that a flat MLP might fail to capture. A hybrid GCN-Transformer architecture, combining spatial graph convolutions with temporal self-attention, represents a particularly promising direction.

Third, temporal multi-scale fusion could combine short-term motion cues (2--3 frames, capturing sudden acceleration and jerk) with medium-term context (8--15 frames, capturing postural transition dynamics) and longer-term behavioral patterns (30+ frames, capturing pre-fall gait instability). A multi-branch architecture with dilated temporal convolutions at different scales, fused via an attention-weighted aggregation, would provide the model with a richer temporal representation than the current single-scale GRU. However, all of these more expressive architectures increase model capacity and thus require substantially more training data to avoid overfitting, further motivating the synthetic data augmentation efforts described in Section~\ref{sec:outlook_synth}.

\subsection{Multi-Camera and Cross-Scene Generalization}
\label{sec:outlook_generalization}

The current model was trained and evaluated exclusively on data from a single Heimann HTPA thermopile sensor array installed at a fixed viewpoint in a controlled indoor environment. While this setup mirrors the intended deployment scenario---a ceiling-mounted sensor in a single room---the resulting model's generalizability to different sensors, viewpoints, and environments remains an open question.

Domain adaptation across different thermal camera models presents a particularly important challenge. Thermal cameras vary widely in spatial resolution (from $32 \times 32$ thermopile arrays to $640 \times 480$ microbolometer arrays), noise equivalent temperature difference (NETD), dynamic range, and spectral sensitivity. A model trained on one camera type may fail catastrophically when deployed on another due to these systematic domain shifts. Unsupervised domain adaptation techniques---such as adversarial feature alignment or optimal transport-based distribution matching---that align the feature distributions of source and target camera domains without requiring labeled data from the target domain will be investigated. Camera-agnostic feature normalization, which removes camera-specific biases in the backbone's intermediate representations, may provide a lightweight alternative.

Multi-view fusion from two or more cameras covering the same physical space opens the possibility of approximate 3D pose reconstruction from 2D thermal inputs. The SensIR project hardware already incorporates two thermopile sensors connected to a single edge device; fusing their overlapping fields of view through stereo calibration and triangulation could yield 3D keypoint coordinates, providing more robust fall detection by reasoning about body orientation relative to the gravity vector rather than relying solely on 2D projections that are inherently ambiguous under viewpoint changes.

Cross-environment generalization---encompassing diverse room layouts, thermal source configurations (radiators, sunlit windows, cooking appliances), furniture arrangements, and population demographics (varying body types, clothing insulation, mobility levels)---will require either collecting labeled data from a broad range of settings or developing self-supervised pre-training strategies. Contrastive learning on large unlabeled thermal video datasets, where the model learns to distinguish temporal neighbors from distant frames and spatial augmentations of the same scene from different scenes, could provide a robust initialization that transfers across environments without extensive per-site fine-tuning.

\subsection{Clinical Integration and Edge Deployment}
\label{sec:outlook_clinical}

Transitioning from offline evaluation on recorded sequences to real-time clinical deployment introduces additional engineering, regulatory, and human-factors challenges that extend beyond pure algorithmic performance.

At the system integration level, the fall detection pipeline must be embedded within a hospital or care facility alert infrastructure. This requires designing a real-time notification workflow that combines instantaneous fall probability with temporal context to minimize nuisance alerts. For example, a clinically meaningful alert policy might require that (i)~the model's fall probability exceeds a threshold for at least two consecutive frames (eliminating single-frame transients), (ii)~the detected person remains in a lying posture for more than 30 seconds without significant movement (distinguishing a fall from a deliberate lie-down), and (iii)~the alert is escalated through a tiered notification system (bedside indicator $\to$ nurse station alert $\to$ emergency page) with configurable delays at each tier. The proposed system's dual-task output---providing both a binary fall score and a four-class posture state---is uniquely suited to this type of composite alerting logic, since posture duration monitoring requires explicit posture classification that single-task approaches do not provide.

On-device deployment targeting the FPGA acceleration hardware specified in the SensIR project, or alternatively NVIDIA Jetson Orin edge GPUs, will necessitate model optimization beyond the current PyTorch implementation. INT8 quantization of the MobileViT backbone and CSPNeXtPAFPN neck, with calibration on a representative thermal IR dataset, is expected to reduce inference latency by approximately $2\times$ with minimal accuracy degradation. Operator fusion, batch normalization folding, and deployment via TensorRT or ONNX Runtime would further reduce overhead. The system's compact 9.3M-parameter footprint and the MobileViT backbone's demonstrated efficiency on mobile hardware~\cite{mehta2022mobilevit} position it well for such optimization. The target throughput of $>$20~FPS on the SensIR FPGA platform appears achievable given the current 80~FPS performance on a desktop GPU.

Beyond acute fall detection, the four-class posture output enables long-term monitoring analytics that are clinically valuable in their own right. By aggregating posture classifications over days and weeks, the system could track daily mobility patterns (time spent standing vs.\ sitting vs.\ lying), sleep regularity (lying onset and wake times), and gradual gait deterioration (decreasing walking episodes or increasing sitting duration) as early predictive indicators of elevated fall risk. These longitudinal trends, presented through a clinical dashboard, could enable proactive interventions---such as physical therapy referrals or environmental modifications---before a fall occurs. Finally, a user study with clinical nursing staff will be essential to evaluate the system's practical utility: measuring alert response times, trust calibration (do staff act on alerts appropriately?), false alarm tolerance thresholds, and the value of posture classification information in daily care workflows.

\subsection{Model Architecture Improvements}
\label{sec:outlook_architecture}

Several architectural refinements could improve the system's accuracy and robustness, particularly as more training data becomes available and the overfitting constraints of the current dataset are relaxed.

The current action head consumes raw keypoint features as an unstructured 35-dimensional vector, without explicitly encoding biomechanical constraints or physical laws governing human movement. Incorporating physics-informed priors could provide strong inductive biases for distinguishing controlled movements from genuine loss-of-balance events. Specifically, center-of-mass (CoM) trajectory estimation---approximated from the seven keypoints using a weighted average with anthropometric segment mass ratios---would provide a direct measure of postural stability. A person whose CoM remains within the support polygon defined by their feet is biomechanically stable; deviation beyond this polygon indicates loss of balance and imminent falling. Joint angle limits (e.g., the knee cannot hyperextend beyond $180^{\circ}$) could serve as soft constraints during pose estimation, improving keypoint accuracy for poses near the boundary of physically plausible configurations. Jerk (the derivative of acceleration) of the head keypoint is a particularly strong fall indicator, as falls exhibit characteristically high jerk values compared to voluntary movements.

The current system employs hand-crafted per-keypoint importance weights that were determined through manual experimentation. Replacing these with learned attention-based keypoint weighting---where a lightweight attention module dynamically weights keypoint contributions based on the current temporal context---would allow the model to discover which keypoints are most informative for each specific posture and action pattern, and to adapt gracefully to occlusion scenarios where certain keypoints are unreliable. For example, the model might learn to rely more heavily on the head and hip keypoints for detecting forward falls but shift attention to the feet during tripping events.

The current two-stage training strategy, while offering modularity and ease of development, introduces an information bottleneck: the pose head is optimized without knowledge of the downstream action classification task, and the action head receives fixed (frozen) keypoint features that may not be optimally informative for fall detection. End-to-end joint training of all three tasks---detection, pose estimation, and temporal action classification---in a single optimization loop could improve gradient flow and feature sharing. The challenge lies in balancing three loss functions with potentially conflicting gradients; techniques such as gradient normalization or uncertainty-based loss weighting could adaptively scale each task's contribution during training.

Finally, the current system classifies each tracked person independently, without modeling spatial relationships between individuals or between persons and environmental objects. In shared care facility spaces, falls frequently involve interaction with furniture (tripping over a chair, falling against a bed railing) or with other residents (collision, assistance attempts). Incorporating a relational reasoning module---such as a scene graph network that models pairwise spatial relationships between detected persons and segmented furniture regions---could improve robustness in these multi-agent scenarios. This extension would require object detection beyond the ``human'' class (adding furniture categories) and a graph neural network for relational inference, representing a substantial but promising architectural expansion.